% chapters/08/evaluasi.tex
\bagian{Asesmen dan Instrumen Penilaian}
\label{sec:asesmen-chase}

\begin{asesmenbox}
\textbf{Kedudukan asesmen bab ini.} Bab ini menuntaskan \SUBCPMK{2.1} dan
merupakan bab pertama pada \CPMK{2}. Ketiga indikatornya diperiksa melalui
penilaian artefak program (Tugas Praktikum 06) dan soal konsep.

\vspace{0.4em}
Materi Bab~VIII termasuk bahan \textbf{Kuis 2} pada minggu kedelapan
(Materi 6--7) dan bahan \textbf{Ujian Tengah Semester} pada minggu kesepuluh
yang mengukur \CPMK{1} dan \CPMK{2}.
\end{asesmenbox}

\subsection*{Instrumen 1: Rubrik Praktikum 06}

\begin{table}[htbp]
\centering
\caption{Rubrik penilaian Praktikum 06 beserta indikator yang diukur}
\label{tab:rubrik-praktikum-06}
{\setlength{\tabcolsep}{4pt}\footnotesize%
\begin{tabularx}{\textwidth}{@{}L{3.2cm} C{1.3cm} C{2.1cm} Y@{}}
\toprule
\headrow \textbf{Kriteria} & \textbf{Bobot} & \textbf{Indikator} & \textbf{Deskripsi capaian} \\
\midrule
Perhitungan arah dan jarak & 35\% & \ind{a} &
Vektor menuju pemain dihitung dengan urutan yang benar, jaraknya dipakai untuk
memutuskan, dan arahnya dinormalkan sebelum dikalikan kecepatan. \\
\addlinespace
Konsistensi kecepatan & 30\% & \ind{b} &
Kecepatan NPC tetap pada segala jarak; tabel pengukuran cara yang benar dan
cara yang salah terisi lengkap beserta penjelasan perbedaannya. \\
\addlinespace
Penjelasan perilaku dan kelemahannya & 35\% & \ind{c} &
Perbedaan AI reaktif dan berbasis perencanaan dijelaskan dengan tepat, dan
catatan praktikum memuat pengamatan sendiri atas kelemahan perilaku reaktif
ketika ada rintangan. \\
\bottomrule
\end{tabularx}}
\end{table}

\subsection*{Instrumen 2: Soal Konsep}

Jawablah dengan singkat. Kunci jawaban beserta indikator yang diukur disajikan
pada Tabel~\ref{tab:kunci-konsep-ch08}.

\begin{enumerate}
  \item Sebuah NPC berada pada $(150, 200)$ dan pemain pada $(450, 600)$.
        Hitunglah vektor menuju pemain, jaraknya, dan vektor kecepatannya
        apabila kecepatannya ditetapkan 100 piksel per detik.
  \item Mengapa vektor arah perlu dinormalkan sebelum dikalikan kecepatan?
        Sertakan satu contoh akibatnya apabila langkah itu dilupakan.
  \item Apa keperluan jarak henti, dan bagaimana Anda menentukan nilainya?
  \item Jelaskan perbedaan AI reaktif dan AI berbasis perencanaan beserta satu
        keadaan permainan yang cocok untuk masing-masing.
  \item Sebuah NPC reaktif menekan dinding selama pemain berada di baliknya.
        Apakah itu kesalahan NPC? Jelaskan.
\end{enumerate}

\begin{table}[htbp]
\centering
\caption{Kunci soal konsep Bab VIII beserta indikator yang diukur}
\label{tab:kunci-konsep-ch08}
{\setlength{\tabcolsep}{4pt}\footnotesize%
\begin{tabularx}{\textwidth}{@{}C{0.8cm} C{2.1cm} Y@{}}
\toprule
\headrow \textbf{No} & \textbf{Indikator} & \textbf{Pokok jawaban} \\
\midrule
1 & \ind{a} &
Vektor menuju pemain $(300, 400)$; jaraknya $\sqrt{300^2 + 400^2} = 500$
piksel; arahnya $(0{,}6, 0{,}8)$; kecepatannya
$(60, 80)$ piksel per detik, dengan panjang 100 piksel per detik. \\
\addlinespace
2 & \ind{b} &
Karena panjang vektor jarak sama dengan jaraknya, sehingga kecepatan menjadi
bergantung pada jarak. Contohnya, pada jarak 300 piksel kecepatannya menjadi
300 piksel per detik, yaitu tiga kali nilai yang ditetapkan pada soal di atas. \\
\addlinespace
3 & \ind{b} &
Agar NPC berhenti ketika sudah cukup dekat, sehingga tidak bergetar di atas
pemain. Nilainya harus lebih besar daripada perpindahan satu bingkai, yaitu
kecepatan dibagi laju bingkai. \\
\addlinespace
4 & \ind{c} &
AI reaktif hanya memakai keadaan saat itu sehingga murah dan mudah diperiksa,
cocok untuk ruang terbuka. AI berbasis perencanaan memperhitungkan rintangan
dan mencari rute, cocok untuk labirin atau ruang dengan banyak rintangan,
dengan biaya perhitungan yang lebih besar. \\
\addlinespace
5 & \ind{c} &
Bukan. NPC itu melaksanakan keputusannya dengan benar berdasarkan keterangan
yang tersedia baginya, yaitu arah menuju pemain. Yang kurang adalah keterangan
mengenai rintangan, dan menambahkannya menuntut pencarian jalur. \\
\bottomrule
\end{tabularx}}
\end{table}
